\documentclass[11pt]{article}

\usepackage[margin=1in]{geometry}
\usepackage{amsmath,amssymb,bm}
\usepackage{booktabs}
\usepackage{listings}
\usepackage{xcolor}
\usepackage{hyperref}

\lstset{
    basicstyle=\ttfamily\small,
    keywordstyle=\color{blue!70!black}\bfseries,
    commentstyle=\color{gray}\itshape,
    stringstyle=\color{green!50!black},
    showstringspaces=false,
    breaklines=true,
    numbers=left,
    numberstyle=\tiny\color{gray},
    frame=single,
    language=Matlab,
}

\newcommand{\norm}[1]{\left\|#1\right\|}
\newcommand{\bv}[1]{\bm{#1}}

\title{Convex Powered--Descent Guidance:\\
       Problem~4 (Acikmese \& Ploen) in MATLAB/CVX}
\author{Companion note for \texttt{p4.m}}
\date{}

\begin{document}
\maketitle

\begin{abstract}
This note documents the optimization problem solved by the script
\texttt{p4.m} and explains how the MATLAB/CVX code realizes it.
The script is a 2D implementation of \emph{Problem~4} from the
enhanced convex formulation of the Mars powered--descent guidance
problem by Acikmese and co--authors~\cite{Acikmese2007,Acikmese2008}.
Given fixed initial and terminal states and a fixed total flight time
$t_f$, the program returns a fuel--optimal thrust profile, formulated
as a second--order cone program (SOCP) through a lossless change of
variables and a Taylor linearization of the mass--flow bounds.
\end{abstract}

% ---------------------------------------------------------------
\section{Problem setting}
% ---------------------------------------------------------------

A lander with dry mass $m_{\text{dry}}$ and initial wet mass
$m_{\text{wet}}$ must be steered from an initial state
$(\bv r_0,\bv v_0)$ to a desired terminal state $(\bv r_f,\bv v_f)$
in a fixed flight time $t_f$, while minimizing fuel.
The planet is Mars, modeled with a constant gravity
vector $\bv g$.  The vehicle has $n_{\text{eng}}$ identical engines
canted by an angle $\phi$ from the body axis; throttle is bounded
between $\eta_{\min}$ and $\eta_{\max}$, and specific impulse
$I_{sp}$ is constant.

The numeric values used in \texttt{p4.m} are listed in
Table~\ref{tab:params}.

\begin{table}[h]
\centering
\begin{tabular}{llr}
\toprule
Symbol & Meaning & Value \\
\midrule
$m_{\text{dry}}$ & dry mass                         & $1505$\,kg \\
$m_{\text{wet}}$ & initial wet mass                 & $1905$\,kg \\
$I_{sp}$         & specific impulse                 & $225$\,s \\
$g_0$            & Earth gravity constant           & $9.80665$\,m/s$^2$ \\
$\bv g$          & Mars gravity vector              & $(0,\,-3.7114)$\,m/s$^2$ \\
$T_{\max}$       & max total thrust at full throttle& $6\times 3100$\,N \\
$\eta_{\min},\eta_{\max}$ & throttle limits          & $0.3,\ 0.8$ \\
$\phi$           & thruster cant angle              & $27^\circ$ \\
$\bv r_0$        & initial position                 & $(2000,\ 1500)$\,m \\
$\bv v_0$        & initial velocity                 & $(100,\ -75)$\,m/s \\
$\bv r_f,\bv v_f$& terminal state                   & $\bv 0,\ \bv 0$ \\
$t_f$            & total flight time                & $75$\,s \\
$\Delta t$       & discretization step              & $1$\,s \\
$N$              & number of knot points            & $t_f/\Delta t + 1 = 76$ \\
\bottomrule
\end{tabular}
\caption{Parameters used in \texttt{p4.m}.}
\label{tab:params}
\end{table}

% ---------------------------------------------------------------
\section{Continuous--time dynamics and lossless convexification}
% ---------------------------------------------------------------

\subsection{Original (non--convex) problem}

Let $\bv r(t)\in\mathbb{R}^2$ be position, $\bv v(t)=\dot{\bv r}(t)$
the velocity, $m(t)$ the mass, and $\bv T(t)\in\mathbb{R}^2$ the
thrust vector.  The translational dynamics are
\begin{align}
\dot{\bv r} &= \bv v, \\
m\,\dot{\bv v} &= \bv T + m\,\bv g, \\
\dot m &= -\alpha\,\norm{\bv T}, &
\alpha &:= \frac{1}{I_{sp}\,g_0\,\cos\phi}.
\label{eq:alpha}
\end{align}
The thrust magnitude is bounded by effective min/max values obtained
by projecting the canted engines onto the body axis,
\begin{equation}
\rho_1 := \eta_{\min}\,T_{\max}\cos\phi,
\qquad
\rho_2 := \eta_{\max}\,T_{\max}\cos\phi,
\qquad
\rho_1 \le \norm{\bv T(t)} \le \rho_2 .
\label{eq:thrust_bounds}
\end{equation}
Fuel--minimization is equivalent to minimizing
$\int_0^{t_f}\norm{\bv T}\,\mathrm{d}t$, or equivalently
maximizing $m(t_f)$.

This problem is non--convex because of the \emph{lower} bound
$\norm{\bv T}\ge\rho_1$ (a non--convex annulus) and the division by
$m$ in the velocity dynamics.

\subsection{Change of variables}

Following Acikmese \& Ploen, define
\begin{equation}
z(t) := \ln m(t),\qquad
\bv u(t) := \frac{\bv T(t)}{m(t)},\qquad
s(t) := \frac{\Gamma(t)}{m(t)},\qquad
\norm{\bv T(t)}\le\Gamma(t),
\label{eq:cov}
\end{equation}
where $\Gamma$ is a scalar slack on thrust magnitude.  In the new
variables the dynamics become
\begin{align}
\dot{\bv r} &= \bv v, \\
\dot{\bv v} &= \bv u + \bv g, \\
\dot z      &= -\alpha\,s, \\
\norm{\bv u} &\le s .
\label{eq:soc}
\end{align}
Relation~\eqref{eq:soc} is a \emph{second--order cone} constraint
and is convex.  Acikmese \& Ploen show that the relaxation
$\norm{\bv T}\le\Gamma$ is lossless: at the optimum it holds with
equality, so no fuel is wasted by working with $\Gamma$ instead of
$\norm{\bv T}$.  The lower thrust bound is mapped to a bound on $s$
that depends on mass:
\begin{equation}
\frac{\rho_1}{m(t)} \le s(t) \le \frac{\rho_2}{m(t)} .
\label{eq:s_bounds_exact}
\end{equation}
Equation~\eqref{eq:s_bounds_exact} is still non--convex in $(z,s)$
because $1/m=e^{-z}$ is convex in~$z$.  A conservative Taylor
expansion around a reference $z_0$ convexifies it, as described in
\S\ref{sec:discrete}.  The objective reduces to
\begin{equation}
\max\; z(t_f) \quad\Longleftrightarrow\quad \max\; m(t_f)
\quad\Longleftrightarrow\quad \min\ \text{fuel used}.
\end{equation}

% ---------------------------------------------------------------
\section{Discretized convex program}
\label{sec:discrete}
% ---------------------------------------------------------------

The trajectory is discretized on $N$ equally spaced knot points
$t_i=(i-1)\Delta t$, $i=1,\dots,N$.  Decision variables are
\[
\bv r_i,\bv v_i,\bv u_i\in\mathbb{R}^2,\qquad
z_i,s_i\in\mathbb{R},\qquad i=1,\dots,N.
\]

\paragraph{Objective.}
\begin{equation}
\max\ z_N .
\end{equation}

\paragraph{Boundary conditions.}
\begin{equation}
\bv r_1=\bv r_0,\quad \bv v_1=\bv v_0,\quad z_1=\ln m_{\text{wet}},
\qquad
\bv r_N=\bv r_f,\quad \bv v_N=\bv v_f .
\end{equation}

\paragraph{Dynamics (trapezoidal with Simpson correction).}
For $i=1,\dots,N-1$,
\begin{align}
\bv v_{i+1} &= \bv v_i + \Delta t\,\bv g
               + \tfrac{\Delta t}{2}\bigl(\bv u_i+\bv u_{i+1}\bigr),
\label{eq:disc_v}\\
\bv r_{i+1} &= \bv r_i + \tfrac{\Delta t}{2}\bigl(\bv v_i+\bv v_{i+1}\bigr)
               + \tfrac{\Delta t^2}{12}\bigl(\bv u_{i+1}-\bv u_i\bigr) .
\label{eq:disc_r}
\end{align}

\paragraph{Mass (trapezoidal).}
For $i=1,\dots,N-1$,
\begin{equation}
z_{i+1} = z_i - \tfrac{\alpha\,\Delta t}{2}\bigl(s_i+s_{i+1}\bigr) .
\label{eq:disc_z}
\end{equation}

\paragraph{Thrust magnitude (SOC).}
For $i=1,\dots,N$,
\begin{equation}
\norm{\bv u_i} \le s_i .
\label{eq:disc_soc}
\end{equation}

\paragraph{Taylor bounds on $s$.}
Let $z_{0,i}$ be a reference value of $z$ corresponding to the
\emph{maximum--thrust} mass history and $z_{1,i}$ the
\emph{minimum--thrust} one:
\begin{align}
z_{0,i} &:= \ln\!\bigl(m_{\text{wet}}-\alpha\,\rho_2\,(i-1)\Delta t\bigr),
\\
z_{1,i} &:= \ln\!\bigl(m_{\text{wet}}-\alpha\,\rho_1\,(i-1)\Delta t\bigr),
\\
\mu_{1,i} &:= \frac{\rho_1}{m_{\text{wet}}-\alpha\,\rho_2\,(i-1)\Delta t}
           = \rho_1\,e^{-z_{0,i}},
\\
\mu_{2,i} &:= \frac{\rho_2}{m_{\text{wet}}-\alpha\,\rho_2\,(i-1)\Delta t}
           = \rho_2\,e^{-z_{0,i}} .
\end{align}
The convex lower and linear upper bounds on $s_i$ are
\begin{align}
s_i &\ge \mu_{1,i}\Bigl[1-(z_i-z_{0,i})+\tfrac{1}{2}(z_i-z_{0,i})^2\Bigr],
\label{eq:taylor_lo}\\
s_i &\le \mu_{2,i}\Bigl[1-(z_i-z_{0,i})\Bigr],
\label{eq:taylor_hi}
\end{align}
together with the physical envelope on $z$,
\begin{equation}
z_{0,i}\le z_i \le z_{1,i} .
\label{eq:z_env}
\end{equation}
Inequality~\eqref{eq:taylor_lo} is a conservative under--approximation
of $\rho_1/m$ (quadratic lower bound of a convex function gives a
convex constraint); \eqref{eq:taylor_hi} is an over--approximation of
$\rho_2/m$ by its tangent line at $z_{0,i}$ (a linear upper bound of
a convex function is convex).  Together they keep the feasible set
convex while tightly tracking the true bounds~\eqref{eq:s_bounds_exact}.

\paragraph{Subsurface constraint (active).}
To forbid flight below the landing pad,
\begin{equation}
r_i^{(2)} \ge -1 \qquad (i=1,\dots,N),
\label{eq:subsurface}
\end{equation}
with $r^{(2)}$ the vertical component.  (The $-1$\,m margin is the
value hard--coded in the script.)

% ---------------------------------------------------------------
\section{Optional constraints (disabled in \texttt{p4.m})}
% ---------------------------------------------------------------

Two additional convex constraints are present in the source but
commented out.  They are included in the formulation here for
completeness.

\paragraph{Thrust--pointing cone.}
Restricting the thrust vector to lie within a half--cone of
aperture $\theta$ around the vertical,
\begin{equation}
u_i^{(2)} \ge s_i\,\cos\theta
\qquad (i=1,\dots,N) .
\label{eq:pointing}
\end{equation}

\paragraph{Glide--slope constraint.}
Keeping the vehicle inside a cone of slope $\gamma$ anchored at
the landing site,
\begin{equation}
r_i^{(1)} \le \frac{r_i^{(2)}}{\tan\gamma}
\qquad (i=1,\dots,N) .
\label{eq:glideslope}
\end{equation}

In \texttt{p4.m} these are the lines that would activate them, but
they are commented out in the current run.

% ---------------------------------------------------------------
\section{MATLAB/CVX implementation walkthrough}
% ---------------------------------------------------------------

The entire program lives in a single self--contained script
\texttt{p4.m}.  It has no dependencies on the other MATLAB files in
this repository except for \texttt{plot\_run2D.m}, which is called
at the end only to visualize the results.

\subsection{Parameter block}

Lines 14--35 declare the numeric values listed in
Table~\ref{tab:params} and pre--compute the constants
$\alpha$, $\rho_1$ and $\rho_2$ defined
by~\eqref{eq:alpha}--\eqref{eq:thrust_bounds}:

\begin{lstlisting}
alpha = 1 / (Isp * g0 * cosd(phi));
r1 = min_throttle * T_max * cosd(phi);
r2 = max_throttle * T_max * cosd(phi);
\end{lstlisting}

\subsection{Decision variables and objective}

Inside the CVX block, five trajectory arrays are declared and
$z_N$ is chosen as the objective:

\begin{lstlisting}
cvx_solver SEDUMI
cvx_begin
    variables r(2,N) v(2,N) u(2,N) z(1,N) s(1,N)
    maximize( z(N) )
    subject to
        ...
cvx_end
\end{lstlisting}

\subsection{Constraint loops}

The constraints are split across two \texttt{for}--loops:

\begin{itemize}
\item A \textbf{dynamics loop} over $i=1,\dots,N-1$ encoding
      \eqref{eq:disc_v}, \eqref{eq:disc_r} and \eqref{eq:disc_z}.
\item A \textbf{pointwise loop} over $i=1,\dots,N$ encoding
      the SOC constraint~\eqref{eq:disc_soc}, the Taylor
      bounds~\eqref{eq:taylor_lo}--\eqref{eq:taylor_hi}, and the
      mass envelope~\eqref{eq:z_env}.  Inside this loop the
      reference quantities $z_{0,i}$, $z_{1,i}$, $\mu_{1,i}$ and
      $\mu_{2,i}$ are recomputed every iteration from the current
      index~$i$.
\end{itemize}

The boundary conditions and the subsurface
constraint~\eqref{eq:subsurface} are vectorized one--liners
outside the loops.  The pointing~\eqref{eq:pointing} and
glide--slope~\eqref{eq:glideslope} constraints are present as
commented--out lines.

\subsection{Post--processing}

After \texttt{cvx\_end}, the script recovers the mass profile
$m_i = e^{z_i}$ and hands the result to the plotting helper:

\begin{lstlisting}
m_vals = exp(z);
plot_run2D(tv, r, v, u, m_vals);
\end{lstlisting}

\texttt{plot\_run2D} produces six figures: position vs.\ time,
velocity vs.\ time, the 2D trajectory with thrust quiver, the
per--axis acceleration, the thrust magnitude, and the acceleration
magnitude.

% ---------------------------------------------------------------
\section{How to run}
% ---------------------------------------------------------------

\begin{enumerate}
\item Install \textbf{MATLAB} and the \textbf{CVX}
      modeling toolbox~\cite{cvx} with the SeDuMi solver enabled
      (\texttt{cvx\_solver SEDUMI}).
\item From the repository root, run
\begin{lstlisting}[numbers=none]
>> p4
\end{lstlisting}
\item Expected output: the CVX solver log, followed by six figure
      windows generated by \texttt{plot\_run2D}.  The optimal
      trajectory reproduces Figure~6 of~\cite{Acikmese2007}.
\end{enumerate}

% ---------------------------------------------------------------
\section{Extension to three dimensions}
\label{sec:3d}
% ---------------------------------------------------------------

This section documents how the convex formulation of
\S\ref{sec:discrete} extends to a full three--dimensional landing
problem.  No 3D implementation is shipped in the current
repository; the material below specifies the mathematical form to
be used by a future \texttt{p4\_3D.m}.

\subsection{Coordinate conventions}

Let $\bv r(t),\bv v(t),\bv u(t)\in\mathbb{R}^3$, with the
convention that $\bv e_3$ points vertically upward.  The Mars
gravity vector becomes
\begin{equation}
\bv g = (0,\,0,\,-3.7114)\ \text{m/s}^2 .
\end{equation}
Note that Acikmese \& Ploen~\cite{Acikmese2007} use $\bv e_1$ as
the radial (vertical) axis; this note adopts the engineering
convention $\bv e_3$ instead so that the 3D code can be obtained
from \texttt{p4.m} by promoting the vertical slot consistently
(the 2D script treats \texttt{r(2,:)} as vertical).

\subsection{Continuous--time dynamics}

The change of variables of \S2.2 carries over unchanged; only the
translational state is promoted to $\mathbb{R}^3$, while the
logarithmic mass $z$ and the thrust slack $s$ remain scalar:
\begin{align}
\dot{\bv r} &= \bv v, &
\dot{\bv v} &= \bv u + \bv g, \\
\dot z      &= -\alpha\,s, &
\norm{\bv u} &\le s .
\end{align}
The lossless convexification argument and the Taylor
bounds~\eqref{eq:taylor_lo}--\eqref{eq:taylor_hi} on $s$ are
dimension--independent and apply verbatim.

\subsection{Discretized 3D convex program}

Decision variables on the same $N$ knot points as in
\S\ref{sec:discrete} are
\[
\bv r_i,\bv v_i,\bv u_i\in\mathbb{R}^3,\qquad
z_i,s_i\in\mathbb{R},\qquad i=1,\dots,N.
\]
The objective ($\max\ z_N$), the boundary conditions, the
trapezoidal dynamics~\eqref{eq:disc_v}--\eqref{eq:disc_r}, the
mass update~\eqref{eq:disc_z}, the Taylor
bounds~\eqref{eq:taylor_lo}--\eqref{eq:taylor_hi}, and the mass
envelope~\eqref{eq:z_env} are identical in form; only the
vector variables now carry a third component.

\paragraph{Thrust magnitude (SOC, $\mathbb{R}^3$).}
\begin{equation}
\norm{\bv u_i}\le s_i\qquad(i=1,\dots,N),
\end{equation}
a 3D second--order cone directly supported by CVX/SeDuMi.

\paragraph{Subsurface constraint.}
\begin{equation}
r_i^{(3)} \ge -1\qquad(i=1,\dots,N) .
\end{equation}

\paragraph{Thrust--pointing cone.}
A half--cone of aperture $\theta$ around the vertical becomes
\begin{equation}
u_i^{(3)} \ge s_i\,\cos\theta\qquad(i=1,\dots,N).
\end{equation}

\paragraph{Axisymmetric glide--slope cone (SOC).}
Keeping the vehicle inside an axisymmetric cone of slope $\gamma$
anchored at the landing site generalizes~\eqref{eq:glideslope} to
\begin{equation}
\sqrt{\bigl(r_i^{(1)}\bigr)^2+\bigl(r_i^{(2)}\bigr)^2}
  \;\le\; \frac{r_i^{(3)}}{\tan\gamma}
\qquad(i=1,\dots,N),
\label{eq:glideslope3d}
\end{equation}
which is a second--order cone constraint and thus convex.

\subsection{Example 3D parameters}

Table~\ref{tab:params3d} lists a suggested 3D initial/terminal
state; the remaining vehicle and discretization parameters are
inherited from Table~\ref{tab:params} to permit direct comparison
with the 2D run.

\begin{table}[h]
\centering
\begin{tabular}{llr}
\toprule
Symbol & Meaning & Value \\
\midrule
$\bv r_0$        & initial position      & $(2000,\ 500,\ 1500)$\,m \\
$\bv v_0$        & initial velocity      & $(100,\ 50,\ -75)$\,m/s \\
$\bv r_f,\bv v_f$& terminal state        & $\bv 0,\ \bv 0$ \\
$\bv g$          & Mars gravity vector   & $(0,\ 0,\ -3.7114)$\,m/s$^2$ \\
$\theta$         & pointing half--angle  & $50^\circ$ \\
$\gamma$         & glide--slope angle    & $4^\circ$ \\
\bottomrule
\end{tabular}
\caption{Suggested 3D boundary conditions and cone angles; all
other parameters ($m_{\text{dry}},m_{\text{wet}},I_{sp},T_{\max},
\eta_{\min},\eta_{\max},\phi,t_f,\Delta t,N$) are taken from
Table~\ref{tab:params}.  These values are indicative and are not
produced by any script in the current repository.}
\label{tab:params3d}
\end{table}

\subsection{Remarks}

Because the lossless convexification structure is
dimension--independent, the 3D program is still a standard SOCP
and can be solved by SeDuMi with only a modest increase in
problem size ($3N$ additional primal variables versus the 2D
case).  Visualization of the 3D trajectory is out of scope here
and would require a companion script \texttt{plot\_run3D.m}
replacing \texttt{plot\_run2D.m}.

% ---------------------------------------------------------------
\begin{thebibliography}{9}

\bibitem{Acikmese2007}
B.~Acikmese and S.~R.~Ploen.
\newblock Convex programming approach to powered descent guidance
         for Mars landing.
\newblock \emph{Journal of Guidance, Control, and Dynamics},
         30(5):1353--1366, 2007.

\bibitem{Acikmese2008}
B.~Acikmese, L.~Blackmore, D.~P.~Scharf, and A.~Wolf.
\newblock Enhancements on the convex programming based powered
         descent guidance algorithm for Mars landing.
\newblock In \emph{AIAA/AAS Astrodynamics Specialist Conference},
         2008.

\bibitem{cvx}
M.~Grant and S.~Boyd.
\newblock CVX: MATLAB software for disciplined convex programming.
\newblock \url{http://cvxr.com/cvx/}.

\end{thebibliography}

\end{document}
